\caption{(a) Convolutional neural network (CNN) architecture used for pose estimation and full-body reconstruction. The $n$ grayscale images from the embedded cameras ($n \times 480 \times 640$) are stacked as input channels and processed by a MobileNetV2 backbone with inverted residual bottlenecks, fine-tuned on our dataset and followed by global average pooling (GAP) to a 1280-dimensional feature vector. Regression heads map this vector to the section-tip poses ($f_p^n$: $6n$ outputs) or to the curvature parameters of each section ($f_C^n$: 5, 4, or 3 parameters per section for the quadratic, affine, and constant models). The heads are one linear layer for $n = 1$ and two layers ($1280 \to 512 \to$ outputs) for $n = 3$. Data augmentation is applied during training only. Transfer learning is applied between pose and curvature networks and from the single-section to the multi-section network. (b) Residual multi-layer perceptron (MLP) used for $N_l^1$ (tip position, 3 $\to$ tendon lengths, 3) and $N_C$ (section-tip pose, 6 $\to$ curvature parameters of the section). The input vector is projected into a 128-dimensional latent space and processed through three residual blocks, each comprising two fully connected layers with Batch Normalization and LeakyReLU activations, followed by the output regression head. The multi-section kinematic networks $N_p^3$ (tendon lengths, 9 $\to$ tip pose, 6) and $N_l^3$ (tip pose, 6 $\to$ tendon lengths, 9) are MLPs with four hidden layers of 128 units and ELU activations; $N_l^3$ is trained through the frozen $N_p^3$ (Eq.~11).}
